\SourceLeaf{355}%
\par\indent Bei einer konformen Abbildung auf die Kugeloberfläche wird ein Dreiecksnetz auf dem Sphäroid, dessen Seiten sogenannte geodätische Linien sind, durch ein entsprechendes Dreiecksnetz dargestellt. Wie schon aus dem Begriff der Konformität unmittelbar folgt, sind dessen Winkel den entsprechenden Winkeln des ursprünglichen Netzes \textit{genau} gleich. Die Seiten werden im strengen mathematischen Sinn zwar keine Großkreisbögen, weichen aber nur sehr wenig von ihnen ab. Lassen sich diese Abweichungen im ganzen Netz so klein halten, dass sie bei der geforderten Rechengenauigkeit als vernachlässigbar gelten dürfen, so ist klar, dass man nach der Übertragung einer Seite des sphäroidischen Netzes auf die Kugeloberfläche das ganze Netz ohne Weiteres wie ein Netz aus gewöhnlichen sphärischen Dreiecken berechnen kann. Erst am Schluss braucht man von den Längen und Breiten auf der Kugeloberfläche zu denen auf dem Sphäroid zurückzugehen, sofern die Endergebnisse der Messung in dieser Form verlangt werden. Dieser Übergang ist entweder mit den Formeln der gewählten Abbildungsart oder mithilfe einer im Voraus berechneten Hilfstafel möglich. In den Fällen dagegen, in denen die Abweichung groß genug ist, um berücksichtigt zu werden, braucht jeder aus den Messungen gewonnene Winkel vor der genauen Rechnung auf der Kugel zunächst eine kleine Reduktion. Das erhöht den Aufwand nur geringfügig, wenn sich die Zahlenwerte der Reduktionen leicht berechnen lassen.

Bei der in der vorliegenden Abhandlung entwickelten Abbildung kann die Abweichung der Kurve, die einen geodätischen Bogen auf der Kugeloberfläche darstellt, vom Großkreisbogen zwischen denselben Endpunkten in der Nähe eines bestimmten Parallelkreises (des Normalparallelkreises) stets als vernachlässigbar gelten. Diesen Parallelkreis kann man frei wählen. Richtet man die gesamte Rechnung für ein bestimmtes Dreiecksnetz selbst von Anfang an ein, so legt man ihn am besten etwa durch die Mitte des ganzen Netzes. Je weiter man sich von diesem Normalparallelkreis nach Norden oder Süden entfernt. desto größer können diese Abweichungen werden. die außerdem sowohl von der Länge der Dreiecksseiten als auch von ihrer Lage zum Meridian abhängen. Selbst in sehr beträchtlicher Entfernung vom Normalparallelkreis bleiben sie jedoch stets so %
